
\begin{proposition}[The weighted-root criterion for elliptic and Hopf trees]\label{cdp:st:C.2.UN.63}
Specify a connected SNC surface with embedded smooth components,
edges between distinct components, and positive component
multiplicities \(m_i\). Each component is a primary or secondary
Hopf surface, or is supplied with its actual fixed-Jacobian,
identity-integral-monodromy elliptic quasi-bundle over a smooth
compact curve, with connected fibres and only smooth multiple
supports. At a supplied quasi-bundle every marked double curve
is a fibre support. Retain every component, analytic torsor and
elliptic gluing map.

Specify the actual weighted line data
\[
 \begin{gathered}
 \sigma_i\in H^0(S_i,\mathcal L_i^{-m_i}),\qquad
 \operatorname{div}\sigma_i=\sum_{j\sim i}m_jJ_{ij},\\
 \mathcal L_i|_{J_{ij}}\simeq
       \phi_{ij}^{*}N_{J_{ji}/S_j},\qquad
 \phi_{ij}^{*}\mathcal L_j|_{J_{ji}}
       \simeq N_{J_{ij}/S_i},
 \end{gathered}
 \tag{C.WT.1}\label{cdp:eq:C.WT.1}
\]
with the equal leading-coefficient line maps at each edge.
Form the finite normalizations
\[
 \mathcal E_i=\operatorname{Norm}
       \{\xi\in\mathcal L_i^{-1}:\xi^{m_i}=\sigma_i\},
 \qquad N=\operatorname{lcm}_i m_i,
 \tag{C.WT.2}\label{cdp:eq:C.WT.2}
\]
their actual elliptic boundary maps and their cyclic
\(\mu_N\)-action. At an edge of multiplicities \(a=gA,b=gB\),
where \(\gcd(A,B)=1\), retain the literal normalized neck
\[
       uv=s^h,\qquad x=u^B,\quad y=v^A,\qquad
       h=\frac{N}{gAB},\qquad t=s^N=x^a y^b.
 \tag{C.WT.3}\label{cdp:eq:C.WT.3}
\]
These are compressed necks; no unspecified contraction of a
reconstructed ruled chain is part of the test.

The following test is necessary and sufficient for the given
weighted support to occur as a fibre of a smooth-total
degeneration of connected complex two-tori.

\begin{enumerate}
\item The root graph is connected. The branch containing an
algebraic-dimension-zero diagonal Hopf component is the complete
two-axis cycle and passes \ref{cdp:st:C.2.UN.59}. In every remaining branch,
each strict elliptic root cover has base genus at most one, and
at most one such cover has genus one.
\item If all elliptic root bases have genus zero, the compressed
root graph is unicyclic. If there is a genus-one root component,
the graph is a tree rooted there. All other root components are
the supplied genus-zero elliptic quasi-bundles or nonlinear
Hopf leaves. There are no unmarked multiple elliptic fibres.
The actual normal lines at each root edge are dual.
\item Let \(m_e\) be their common exact torsion order, and let
\(h_e\) be the exponent in \eqref{cdp:eq:C.WT.3}. Give the cycle, or the
genus-one root, height zero, and set
\[
       H_v=\sum_{e\text{ on the rooted path to }v}h_e m_e.
 \tag{C.WT.4}\label{cdp:eq:C.WT.4}
\]
All nonlinear leaves have one common height \(R>0\);
every elliptic vertex has height strictly less than \(R\).
When there is no nonlinear leaf, \(R\) may be chosen larger
than all the heights.
\item Perform the exact inverse operations below in whole
cyclic-group orbits, retaining the complete torsors, the full
secondary integers, primitive fibre bases and all marked maps.
The final cycle has actual rank-three complement product
\(M=I\), or the final principal elliptic torsor over its
elliptic base has integral Chern vector zero.
\item Let \(D\) be the finite linear part on this actual pruned
core. For a cycle, it is the linear part on its common
two-dimensional logarithmic vector space; for a torus core,
it is the linear part of its actual affine torus action.
For the root generator with base action \(s\mapsto\zeta_Ns\),
the height satisfies
\[
                         \zeta_N^R=\det_{\C}D.
 \tag{C.WT.5}\label{cdp:eq:C.WT.5}
\]
In the absence of nonlinear leaves one chooses any sufficiently
large positive \(R\) in this residue class.
\end{enumerate}

Every step uses finite covers, actual line and torsor maps,
integral lattices, and the displayed inverse charts. In
particular \eqref{cdp:eq:C.WT.5} is an equality of the actual canonical
character. A cycle action may reverse orientation through
vertices; such reflections are included.

Every support passing the test has a realization over
\(\PP^1\) with smooth compact total threefold. The construction
retains its original components, gluing maps and all
multiplicities. It can have exactly the prescribed weighted
support at zero and its complex conjugate at infinity, with
all other fibres smooth connected two-tori. Every reduced
fibre support has Euler number zero. Normal-product constants
are chosen in the realization; no equivalence of different
full analytic neighbourhoods is asserted.
\end{proposition}

\begin{proof}
\par\medskip\noindent\textbf{The actual component covers and their local inertia.}
An actual map supplies \eqref{cdp:eq:C.WT.1} by dividing its defining
function by the \(m_i\)-th power of the canonical divisor
section and restricting to \(S_i\). On its strict root
component the exact formula is
\(\xi=s^{N/m_i}/v_i\), where \(v_i\) is that normal
divisor coordinate. Thus the root generator is
\[
                  \xi\longmapsto
                         \zeta_N^{N/m_i}\xi.
 \tag{C.WT.6}\label{cdp:eq:C.WT.6}
\]
At a double curve, absorbing a unit in one normal
coordinate gives \(x^a y^b=s^N\). Its \(g\) local
branches have precisely the finite birational
normal maps \eqref{cdp:eq:C.WT.3}, as proved in (C.WH.4).
The equal leading-coefficient maps in \eqref{cdp:eq:C.WT.1}
identify their actual curve covers and normal
characters. This constructs the compressed root
data without replacing a nonlinear Hopf
neighbourhood by its normal line.

Every strict component cover is smooth. Locally along
a marked curve its equation is
\(\xi^m=z^b\) times a unit. Write
\(d=\gcd(m,b)\), \(q=m/d\), \(B=b/d\).
On each branch its normalization is
\[
                         z=w^q,\qquad \xi=cw^B,
 \tag{C.WT.7}\label{cdp:eq:C.WT.7}
\]
with \(c\) a fixed root of unity. This is a smooth
coordinate chart with the unchanged curve
coordinate. Away from these curves the map is
finite \'{e}tale. The induced finite maps of
reduced boundary curves are therefore \'{e}tale.
Every stabilizer preserving such a connected
elliptic curve acts on it by deck translations.
The same is true on its connected normal-root
elliptic cover.

For a supplied elliptic component, remove the finitely
many marked base values. A connected component of
its finite surface cover has a fixed finite-index
fibre subgroup \(\Lambda'\subset\Lambda_F\).
Take the actual connected-fibre Stein base cover.
On it the fibre is the fixed elliptic curve
\(F'=\C/\Lambda'\), with the actual isogeny
\(\pi:F'\to F\). A lifted transition \(a\)
satisfies
\[
                            \pi a=T_b\pi.
 \tag{C.WT.8}\label{cdp:eq:C.WT.8}
\]
Its derivative is therefore the identity, so it
is a translation. This proves identity integral
monodromy on the connected-fibre base cover,
including loops which permuted the original
components of a fibre cover.

Near a deleted value lift the actual chart
\((F\times\Delta)/\langle(\tau,\zeta_m)\rangle\)
and apply \eqref{cdp:eq:C.WT.7} to its normal power. Quotient
first by the subgroup trivial on \(F'\); it
acts by normal pseudoreflections. The remaining
inertia acts freely by translations on \(F'\).
The result is again a smooth multiple-fibre
chart
\((F'\times\Delta)/\langle(\tau',\zeta_{m'})\rangle\).
The finite translation gauge of (C.HT.4) makes
\(\tau'\) constant. Thus the compact strict cover
has an actual fixed-Jacobian quasi-bundle over
the smooth compact Stein base \(B'\), with only
smooth multiple supports. Its invariant
cohomology frame gives
\[
                  R^1q_*\OO=\OO_{B'},\qquad
                         p_g=g(B').
 \tag{C.WT.9}\label{cdp:eq:C.WT.9}
\]
This argument also applies to every
algebraic-dimension-one Hopf component, using
its actual fibration from \ref{cdp:st:C.2.UN.37}.

For a nonlinear Hopf component there is a
different explicit map. On its primary cover
write the contraction
\[
                    (z,w)\longmapsto
                     (\beta^r z+c w^r,\beta w),
                 \qquad c\ne0 .
 \tag{C.WT.10}\label{cdp:eq:C.WT.10}
\]
Its only curve is \(w=0\). The power of the
weighted normal line is the constant-automorphy
axis line; dividing its transferred connection
gives a flat connection on the line itself.
On \(\C^2\setminus0\) the root section is a
constant times \(w^b\). Indeed its quotient
by \(w^b\) extends as a holomorphic unit at
zero; its automorphy factor is one by
evaluation there, and contraction makes that
unit constant. Formula \eqref{cdp:eq:C.WT.7} is consequently
the actual map \((z,v)\mapsto(z,v^q)\).
A positive power of \eqref{cdp:eq:C.WT.10} preserving one
branch lifts to
\[
                 (z,v)\longmapsto
                    (\alpha^{qr}z+c'v^{qr},\alpha v),
                     \qquad c'\ne0 .
 \tag{C.WT.11}\label{cdp:eq:C.WT.11}
\]
Its branch stabilizer is free on
\(\C^2\setminus0\). A fixed point would project
to a fixed point of the original free deck
group; an element over its identity fixing
both \(v^q\) and \(v^B\) fixes \(v\), since
\(\gcd(q,B)=1\). Hence the finite normalization
is exactly a primary or secondary nonlinear
Hopf surface, with the displayed map and
resonance \(qr\).

The nontorsion axis normals of a diagonal
algebraic-dimension-zero Hopf cannot satisfy
the weighted normal equality with a torsion
normal. Its axis covers remain nonresonant:
a multiplicative relation after taking finite
coordinate roots would give a multiplicative
relation for the original multipliers.
\ref{cdp:st:C.2.UN.40} forces both its axes into the boundary.
Thus that connected branch is the complete
cycle already covered by \ref{cdp:st:C.2.UN.59}. We henceforth
treat the torsion-normal branch.

Resolve the \(A_{h-1}\) necks only to apply
the compact-core inequality \ref{cdp:st:C.2.UN.28}. It gives
\(\sum p_g\le1\) for all strict components.
By \eqref{cdp:eq:C.WT.9}, their base genera are at most one
and at most one is one. No assertion that
a many-marked Hopf root cover is still Hopf
has been made.

\par\medskip\noindent\textbf{Compressed canonical heights and the actual period character.}
The central fibre of \(uv=s^h\) is reduced,
and its abstract two-component crossing is
SNC. Resolving the neck inserts the actual
elliptic ruled chain, with the same normal
root maps. Either description gives the
same component and edge cohomology maps.
The exact restriction square (C.QT.2)--(C.QT.3)
and the normalization and residue arguments
of \ref{cdp:st:C.2.UN.61}--\ref{cdp:st:C.2.UN.62} therefore give the following.
With no genus-one vertex the root graph is
unicyclic, and the minimum canonical order
is attained exactly on its cycle. With one
genus-one vertex the graph is a tree and
that vertex is the unique minimum. One can
also apply \ref{cdp:st:C.2.UN.47} to the latter assertion.
The connectedness here follows from the
proper Stein factor of the actual root
family, whose nearby fibres are connected.

Here is the precise change caused by a
compressed neck. Use a curve coordinate
\(z\) and the relative canonical frame
\(du/u\wedge dz=-dv/v\wedge dz\).
If the leading form on the \(u\)-side has
curve order \(d\), then
\[
 s^k u^{d+1}\frac{du}{u}\wedge dz
      =-s^{k+h(d+1)}v^{-d-1}
                         \frac{dv}{v}\wedge dz .
 \tag{C.WT.12}\label{cdp:eq:C.WT.12}
\]
Thus the canonical-order difference is
\(h(d+1)\), rather than \(d+1\).
At an elliptic multiple fibre of order \(m\),
the exact wedge order is
\(d=m\operatorname{ord}\theta+m-1\).
Consequently the edge integer is
\[
       \frac{k_w-k_v}{h_e m_e}
                        =\operatorname{ord}_{p_e}\theta+1.
 \tag{C.WT.13}\label{cdp:eq:C.WT.13}
\]
At an unmarked multiple fibre the required
order is zero; the formula excludes every
unmarked \(m>1\). The degree of the actual
base differential gives the same
valence-minus-two equation as \ref{cdp:st:C.2.UN.58}.
The finite leaf induction now has outward
increment \(h_e m_e\). This proves exactly
\eqref{cdp:eq:C.WT.4}, including product caps with
normal order one. Nonlinear terminal
forms give the same increment, by \eqref{cdp:eq:C.WT.10}.

The root collar sends the connected cover
\(u^{m_e}=c\) to
\(v^{m_e}=s^{h_e m_e}/c\). These are the
same oriented integral two-cycle. Together
with the full secondary boundary matrix
of \ref{cdp:st:C.2.UN.50}, they identify every regular
elliptic fibre and every nonlinear core
with one integral class \(\gamma_s\).
The canonical period has order greater
than \(H_v\) on an elliptic vertex, and
has order exactly \(H_v\), with the
nonzero coefficient of (C.PT.10), at a
nonlinear leaf. This proves the height
test even on an arbitrary actual
smoothing.

Choose a finite-equivariant relative
canonical generator: a nonconstant unit
in its finite character is removed by
the finite averaging of its logarithm.
Divide by the minimum power of \(s\).
The inverse charts used below carry this
form to the nonzero canonical frame on
the pruned core. Its finite character
there is exactly \(\chi=\det_{\C}D\).
The deck maps preserve the oriented
regular-fibre class, so
\[
 P(\zeta_Ns)=\chi P(s),\qquad
 P(s)=\int_{\gamma_s}\Omega_s .
 \tag{C.WT.14}\label{cdp:eq:C.WT.14}
\]
If a nonlinear leaf is present, its
nonzero leading coefficient proves
\(\zeta_N^R=\chi\). Without such a leaf
one may choose a larger positive exponent
in this character class. This proves
\eqref{cdp:eq:C.WT.5} with the actual minimum-form
character, without inserting an
unproved parity assumption.

\par\medskip\noindent\textbf{Equivariant inverse pruning and the exact core obstruction.}
Use the one-mark root-torsor classification
(C.PR.1)--(C.PR.8), or its identical application
to the supplied quasi-bundles in \ref{cdp:st:C.2.UN.61}--\ref{cdp:st:C.2.UN.62}.
After all children of a vertex have been
removed, its only possible multiple fibre
is its parent mark. Its actual torsor is
therefore a recorded diagonal Hopf pair,
with its full integer \(n\), or the product
ruled cap. At a Hopf step choose
\(mk-nl=1\). The actual inverse maps are
\[
 \begin{gathered}
 (u,X)\sim(\gamma^m u,\gamma^lX),\qquad
 p=X^m u^{-l},\qquad W=X^{-n}u^k,\\
                       u=p^nW^m,\qquad X=p^kW^l.
 \end{gathered}
 \tag{C.WT.15}\label{cdp:eq:C.WT.15}
\]
They identify the punctured parent tube
with the specified product tube. The
parameter \(p\) is its actual base
coordinate, since its pullback is the
displayed multiple-fibre function.
The tensor and torsor maps therefore
extend across the new ordinary fibre,
with all other base charts unchanged.
The identity
\[
             d\log p\wedge d\log W
                         =d\log X\wedge d\log u
 \tag{C.WT.16}\label{cdp:eq:C.WT.16}
\]
also proves the canonical-frame comparison
used in \eqref{cdp:eq:C.WT.14}: the old order \(m-1\)
becomes the nonzero frame \(dp\wedge
d\log W\). The normal-product neck
exponent \(h\) changes the order in the
family parameter, but does not change
this central inverse map.

Carry out these operations simultaneously
on whole finite-group orbits. A stabilizer
of a connected boundary curve acts there
by a deck translation, as proved after
\eqref{cdp:eq:C.WT.7}. On the connected normal-root
elliptic cover its linear part is also
one, by the isogeny derivative argument
of \eqref{cdp:eq:C.WT.8}. It consequently fixes the
entire primitive fibre lattice. Choose
one primitive basis and the full integer
\(n\) on each orbit and transport it.
The maps \eqref{cdp:eq:C.WT.15}, their canonical
divisor sections and their inverses
then give the finite maps on the
replacement tubes. Their group law
holds on the punctured tubes and hence
on their holomorphic extensions.
This retains the stabilizer characters;
no independent choice of roots of
transition multipliers is made.
Disjoint replacements commute by their
actual atlases, so a group-invariant
postorder gives the required exact
equivariant pruning record.

The root action cannot invert an edge:
its two endpoints lie over two distinct
original components, each fixed by the
root generator. It can, however, reflect
a cycle through vertices, interchanging
two different preimage curves over one
original boundary. Tree parent
orientations are preserved.

For necessity no analytic contraction
of the given family is used.
At a fixed \(s\ne0\) the neck
\(uv=s^h\) has the same boundary
homomorphisms as \(uv=s\).
The full secondary matrix in \ref{cdp:st:C.2.UN.50}
shows that removal of the compact
leaf piece kills precisely the
primitive meridian \(\delta\) and
retains the two fibre generators.
The inverse product filling has
those identical maps. The actual
Van Kampen pushouts of (C.PR.9)
therefore preserve the nearby
fundamental group throughout.
For a cycle the final injected
semidirect product is
\(\Lambda\rtimes_M\Z\); its equality
with \(\Z^4\) forces \(M=I\).
For a genus-one core the injected
fibre lattice in (C.PT.12) has
commutator its integral Chern
vector, so that vector is zero.
The explicit Cech lattice
(C.PT.8) identifies this very
principal torsor with a complex
two-torus. This proves the fourth
test while retaining the full
analytic class.

In the cycle case the exact two-mark
classification gives marked diagonal
Hopf or degree-zero ruled components.
Their regular fibres are the common
compact elliptic subgroup of their
open group \(\C^2/\Lambda\).
The finite maps on that group are
affine. For a diagonal Hopf pair,
an axis-preserving lift has form
\((zU,wV)\); contraction makes
the units \(U,V\) constant, exactly
as in \ref{cdp:st:C.2.UN.59}. An axis-interchanging
lift is the composition of such
a map with the coordinate
exchange. Marked ruled maps have
the actual affine logarithmic
lifts of \ref{cdp:st:C.2.UN.51}. Hence the finite
linear part \(D\) is a single
map on the common vector space.
For a torus core its derivative
is constant by periodicity,
so its actual finite action is
affine as well.

\par\medskip\noindent\textbf{Signed cycle equations, including reflections.}
Orient the pruned cycle and let \(n_j\)
be the outgoing normal meridian,
\(h_j\) its retained neck exponent.
Let \(\epsilon=1\) or \(-1\) be the
orientation character of the root
generator, and \(\sigma\) its
permutation of vertices. The image
edge is \(k=\sigma(j)\) for
\(\epsilon=1\), and
\(k=\sigma(j+1)\) for \(\epsilon=-1\).
In actual logarithmic coordinates
write
\[
 \begin{gathered}
 g(j,v,s)=(\sigma j,Dv+\tau_j,\zeta_Ns),\\
 e_j(v,s)=v+b_j+l_j\log s,\qquad
       l_j=-\frac{h_jn_j}{2\pi i},\\
 Dn_j=\epsilon n_k,\qquad
 Dl_j=\epsilon l_k.
 \end{gathered}
 \tag{C.WT.17}\label{cdp:eq:C.WT.17}
\]
For a reverse edge the comparison
is with \(e_k^{-1}\), not \(e_k\).
Taking \(\log\zeta_N=2\pi i/N\),
the actual compositions give
\[
 \begin{aligned}
 \tau_{j+1}+Db_j-\tau_j
       -\epsilon b_k-\epsilon l_k\,2\pi i/N
                                      &\in\Lambda,\\
 \sum_{a=0}^{N-1}D^{N-1-a}
                       \tau_{\sigma^a j}&\in\Lambda.
 \end{aligned}
 \tag{C.WT.18}\label{cdp:eq:C.WT.18}
\]
For example when \(\epsilon=-1\),
\(g_{j+1}e_j=e_k^{-1}g_j\) gives
the first line directly. These
equalities retain every affine
constant of the actual central
action.

Put \(b=\sum b_j\) and
\(u=-\sum h_jn_j=dp\), with \(p\)
primitive and \(d>0\). The exact
contraction-height calculation
(C.AD.3) gives
\(\ell(in_j)\) of one nonzero
sign, where
\(\ker\ell=\operatorname{span}_{\R}\Lambda\).
For ruled vertices this is the
actual elliptic quotient
calculation of \ref{cdp:st:C.2.UN.55}.
Thus \(\ell(iu)\ne0\) and
\(Du=\epsilon u\).
Choose the common elliptic
line and \(\C p\) as coordinates.
Then
\[
 \begin{gathered}
 W=\operatorname{span}_{\R}\Lambda=\C\times\R,\quad
 p=(0,1),\quad D=\operatorname{diag}(\rho,\epsilon),\\
 \Lambda=\langle(a_1,c_1),(a_2,c_2),p\rangle_{\Z},
 \qquad c_1,c_2\in\R,\qquad
        \operatorname{Im}(\overline{a_1}a_2)>0,\\
              (D-\epsilon I)b-\epsilon u/N\in\Lambda.
 \end{gathered}
 \tag{C.WT.19}\label{cdp:eq:C.WT.19}
\]
The last line is the sum of
\eqref{cdp:eq:C.WT.18}. The finite elliptic
linear part \(\rho\) has order
\(1,2,3,4\) or \(6\), by its
integral trace, as in \ref{cdp:st:C.2.UN.59}.
In particular
\(\det D=\rho\epsilon\), including
the sign at a reflected core.

Adding \(c n_j\) to \(b_j\) on
a whole edge orbit, with the
same \(c\in\C\), preserves
\eqref{cdp:eq:C.WT.18}, since
\(Dn_j=\epsilon n_k\).
The second coordinates of all
the \(n_j\) have the same
nonzero sign. Such an orbit
therefore allows the second
coordinate of \(b\) to be
set to \(i\). This changes
only a normal-product constant:
the translation along the
normal meridian restricts
to zero on the double curve.
It preserves the specified
elliptic gluing maps and the
entire finite action.

\par\medskip\noindent\textbf{The finite-equivariant compact core.}
For a cycle use the coordinates
in \eqref{cdp:eq:C.WT.19}; for a torus choose
a finite-invariant Hermitian
metric and an orthonormal
complex basis. In either case
\(D\) is unitary and
\[
 D J\overline D^{-1}=(\det D)J,
       \qquad J=\begin{pmatrix}0&1\\-1&0\end{pmatrix}.
 \tag{C.WT.20}\label{cdp:eq:C.WT.20}
\]
This follows by multiplying
the two-by-two matrices, or
from the alternating form
on a two-dimensional space.
With \(R\) as in \eqref{cdp:eq:C.WT.5} set
\[
 V_s(w)=w+s^R J\overline w,\qquad
             D V_s=V_{\zeta_Ns}D .
 \tag{C.WT.21}\label{cdp:eq:C.WT.21}
\]
For a torus core the four
real period vectors \(V_s\Gamma\)
remain independent because
\(V_s\) has real determinant
\((1+|s|^{2R})^2\).
Transport an affine constant
to the target frame:
\(\tau(s)=V_{\zeta_Ns}\tau\).
All finite lattice and order
equalities are then preserved
exactly, so these are an actual
finite action on the torus
family.

For a cycle transport
\[
 b_j(s)=V_s b_j,\qquad
 \tau_j(s)=V_{\zeta_Ns}\tau_j,\qquad
 B_s=V_s b+\frac{d\log s}{2\pi i}V_s p .
 \tag{C.WT.22}\label{cdp:eq:C.WT.22}
\]
The correction
\(l_j\,2\pi i/N=-h_jn_j/N\)
in \eqref{cdp:eq:C.WT.18} is a real-rational
period vector. Applying \(V_s\)
therefore carries the signed
equations to the identical
equations in the target frame.
Summing gives
\[
                 D B_s-\epsilon B_{\zeta_Ns}
                               \in V_{\zeta_Ns}\Lambda .
 \tag{C.WT.23}\label{cdp:eq:C.WT.23}
\]
The finite action consequently
preserves the actual rank-four
nearby period lattice. The
global map itself can be
written on the actual middle
charts. Choose a cut and put
\[
 \begin{gathered}
 \theta_j(s)=\sum_{a<j}
      \bigl(b_a(s)+l_a(s)\log s\bigr),\qquad
 \eta=v_j-\theta_j(s),\\
 \eta\longmapsto
 D\eta+D\theta_j(s)+\tau_j(s)
                         -\theta_{\sigma j}(\zeta_Ns),
 \qquad l_a(s)=-h_aV_s(n_a)/(2\pi i).
 \end{gathered}
 \tag{C.WT.23a}\label{cdp:eq:C.WT.23a}
\]
Changing an adjacent chart
changes this expression by
the actual lattice element
in \eqref{cdp:eq:C.WT.18}. Changing the
cut adds \(B_s\) on the source
and \(\epsilon B_{\zeta_Ns}\)
on the target, so \eqref{cdp:eq:C.WT.23}
proves that the expression
is unchanged as a map of
the middle torus. A logarithm
change adds integral normal
periods. The second equation
of \eqref{cdp:eq:C.WT.18} proves its
\(N\)-th-order relation.
Thus this is an actual
global finite holomorphic
action, including reflections,
not only an invariant lattice.
The
normal maps are the literal
\(uv=s^{h_j}\) maps; reversing
an edge exchanges their two
coordinates. The moving
Hopf and ruled frame maps
of \ref{cdp:st:C.2.UN.57}--\ref{cdp:st:C.2.UN.59} give these
actual charts also for a
reflection through a vertex.

The core is compact over the
whole parameter sphere. In
the cycle case its determinant
is a positive constant times
\[
 (1+r^{2R})
 \left[1+r^{2R}
    +\frac d{2\pi}(r^{2R}-1)\log r\right],
                \qquad r=|s|>0 .
 \tag{C.WT.24}\label{cdp:eq:C.WT.24}
\]
Each factor is positive.
This is the exact computation
of \ref{cdp:st:C.2.UN.57} with \(V_s\) replaced
by \(V_{s^R}\), keeping
\(d\log s\) unchanged.
It gives compact connected
two-tori on every punctured
fibre. Finite core and neck
collars give the proper
filling at zero. No torus
rank is inferred merely
from the fundamental group.

At infinity put \(z=s^{-1}\)
and divide the vector frame
by \(s^R\). The periods become
\[
              W_z(w)=J\overline w+z^Rw,\qquad
              D_{\infty}=(\det D)^{-1}D
                       =J\overline D J^{-1}.
 \tag{C.WT.25}\label{cdp:eq:C.WT.25}
\]
The base character is
\(\zeta_N^{-1}\). An affine
constant becomes
\(J\overline\tau+(\det D)^{-1}z^R\tau
=W_{\zeta_N^{-1}z}(\tau)\).
The normal meridians are
\(-J\overline{n_j}\); changing
\(\log s\) to \(-\log z\)
gives the conjugate core
with the same positive neck
exponents. Thus the infinity
maps, including every affine
constant and reflection, are
the actual complex conjugates
of the central core maps.
This provides the second
proper compact end.

\par\medskip\noindent\textbf{The precise finite jet gauge at a surgery orbit.}
The common regular elliptic fibre has a
primitive period basis \(\lambda,\mu=c\lambda\),
\(\operatorname{Im}c>0\). In the core frame
its exact determinant is
\[
 \det_{\C}(V_s\lambda,V_s\mu)
              =2i\operatorname{Im}(c)\|\lambda\|^2s^R .
 \tag{C.WT.26}\label{cdp:eq:C.WT.26}
\]
Consequently the actual rank-two period
cover near a marked fibre has the
semi-torus presentation
\[
 (p,w,s)\longmapsto(p+\delta(s),\gamma(s)w,s),
 \qquad \delta(s)=s^{\ell} d(s),\quad d(0)\ne0,
                     \quad|\gamma(0)|<1 .
 \tag{C.WT.27}\label{cdp:eq:C.WT.27}
\]
Initially \(\ell=R\). The exact primitive
basis change (C.PR.13) multiplies \(\delta\)
by a holomorphic unit and retains its
order. Every later chart below has this
same presentation.

Consider the stabilizer of one surgery
curve, of order \(d_0\), with base
character \(\xi\). This is the stabilizer
order, not the order of the whole
graph group when that curve moves.
As proved after \eqref{cdp:eq:C.WT.15}, its linear
action on the actual fibre lattice
is the identity. On the universal
logarithmic cover its finite lift
therefore has an upper-triangular
affine form
\[
 (p,V,s)\longmapsto
    (a(s)p+b(s),\,V+c(s)p+e(s),\,\xi s),
                     \qquad V=\log w .
 \tag{C.WT.28}\label{cdp:eq:C.WT.28}
\]
The lift may have \(d_0\)-th power
equal to the \(k\)-th infinite deck
translation, with an integral
angular deck translation as well.
Both integers are retained from the
actual root action. The \(p\)-part
of that equality is
\[
 \prod_{j=0}^{d_0-1}a(\xi^j s)=1,\qquad
                 g^{d_0}(p)=p+k\delta(s).
 \tag{C.WT.29}\label{cdp:eq:C.WT.29}
\]
The norm-one unit cocycle for \(a\)
is made constant \(a_0=a(0)\) by
a holomorphic unit gauge, equal
to one at zero. Explicitly take
a logarithm of \(a/a_0\), whose
finite orbit sum is zero, and
use the same finite weighted
sum as in (C.HT.4). Its
coboundary is that logarithm.
Commutation with the infinite
period then gives
\[
 \delta(\xi s)=a_0\delta(s),\qquad
 \sum_{j=0}^{d_0-1}a_0^{d_0-1-j}b(\xi^j s)
                                      =k\delta(s).
 \tag{C.WT.30}\label{cdp:eq:C.WT.30}
\]

For any coefficient \(b_j s^j\)
of order \(j<\ell\), the sum
in \eqref{cdp:eq:C.WT.30} is zero unless
\(\xi^j=a_0\); in that resonant
case it is
\(d_0 a_0^{d_0-1}b_j s^j\),
so \(b_j=0\). Every nonresonant
coefficient is removed by
\(p\mapsto p+c_j s^j\), since
its change is
\((\xi^j-a_0)c_j s^j\).
Thus, whenever \(\ell\ge hm\),
one makes
\[
                               b(s)=O(s^{hm})
 \tag{C.WT.31}\label{cdp:eq:C.WT.31}
\]
by an actual finite polynomial
jet gauge. This preserves the
central marked point and every
resonant coefficient. For an
intermediate diagonal Hopf
component a constant nontrivial
affine action is first centered
at its finite fixed point;
this is its actual triangular
Hopf automorphism and is
recorded on all remaining
marks.

In particular the nonlinear
translation is never discarded.
On a nonlinear Hopf universal
cover put \(U=z/w^m,V=\log w\).
Its infinite generator is
\((U,V)\mapsto(U+c,V+\log\beta)\),
with \(c\ne0\). A parent-axis-
preserving lift commuting with
a positive central contraction
has \(U\mapsto U+b\):
an affine coefficient \(a\)
must satisfy \(ac=c\).
If its \(d_0\)-th power is the
\(k\)-th infinite generator
times a finite angular element,
that element fixes \(U\), so
\[
                                d_0b=kc.
 \tag{C.WT.32}\label{cdp:eq:C.WT.32}
\]
This is exactly the term which
survives \eqref{cdp:eq:C.WT.30} at a terminal
insertion. Indeed if \(\ell=hm\),
then \(a_0=\xi^{\ell}\) and the
leading coefficient there is
\(b_{\ell}=k a_0d(0)/d_0\).
After dividing by \((\xi s)^{hm}\)
it is \(k d(0)/d_0\), as required
by \eqref{cdp:eq:C.WT.32}. For a diagonal
component the corresponding
equation is
\(b(1+a+\cdots+a^{d_0-1})=0\);
it forces \(b=0\) when \(a=1\)
and permits exactly the finite
fixed-point gauge otherwise.
This proves the full central
triangular comparison.

\par\medskip\noindent\textbf{The literal compressed weighted insertion.}
For a recorded Hopf step of
normal order \(m\), full secondary
integer \(n\), and neck exponent
\(h\), use the actual normalized
pullback, by \(t=s^h\), of the
weight-\((m,1)\) base modification
in \ref{cdp:st:C.2.UN.60}. Its root chart is
\[
       p=x^m Z,\qquad xW=s^h,\qquad (Z,W)\ne(0,0).
 \tag{C.WT.33}\label{cdp:eq:C.WT.33}
\]
Retain the pullback of the old
tautological line \(\OO(nE_{\rm old})\).
Its canonical local section is
\(x^n\); since \(x=s^h/W\) at a
generic point of the new component,
\(\operatorname{div}(x)=hE_{\rm new}\).
Thus the pulled line has coefficient
\(nh\), not \(n\), on this generic
divisor. Replacing it by
\(\OO(nE_{\rm new})\) would change
the actual root action and is not
performed here.

Choose the recorded \(\beta(s)\)
with \(\beta(s)^n=\gamma(s)\).
The nonzero line has the free
finite chart
\[
       (Z,W,x)\longmapsto
                   (\eta^mZ,\eta W,\eta^{-1}x),
                           \qquad\eta\in\mu_n.
 \tag{C.WT.34}\label{cdp:eq:C.WT.34}
\]
It is free because
\(\gcd(m,n)=1\), including
where one of \(Z,W\) vanishes.
The infinite deck map is exactly
\[
 \begin{split}
 (Z,W,x,s)\longmapsto\bigl(
   &\beta(s)^m[Z+d(s)s^{\ell-hm}W^m],\\
   &\beta(s)W,\ \beta(s)^{-1}x,\ s\bigr).
 \end{split}
 \tag{C.WT.35}\label{cdp:eq:C.WT.35}
\]
It preserves \(xW=s^h\) and
changes \(p=x^mZ\) by precisely
\(d(s)s^{\ell}=\delta(s)\).
Every term has the same finite
character as \(Z\) in
\eqref{cdp:eq:C.WT.34}. The central child
is nonlinear exactly when
\(\ell=hm\), and diagonal
when \(\ell>hm\). Its recorded
nonzero nonlinear coefficient
is matched by the actual
rescaling of \(Z\); the finite
translation in \eqref{cdp:eq:C.WT.32}
scales with it.

Away from its parent curve put
\(U=Z/W^m\) and use the finite
coarse fibre coordinate \(W^n\).
The exact new deck is
\[
            (U,W^n)\longmapsto
            (U+\delta(s)/s^{hm},\,\gamma(s)W^n).
 \tag{C.WT.36}\label{cdp:eq:C.WT.36}
\]
This is the required remaining
order \(\ell-hm\), at every
other marked child fibre.
There is no change of the
given analytic torsor outside
the one recorded tube.

The finite action also lifts
exactly. The gauge \eqref{cdp:eq:C.WT.31}
makes its base map preserve
the weighted ideal
\((p,s^{hm})\), and hence its
normalized weighted modification.
The canonical divisor section
gives the action on the pulled
tautological line. In a chart,
an \(n\)-th root of a nonzero
multiplier may be chosen;
different choices are precisely
the \(\mu_n\)-action \eqref{cdp:eq:C.WT.34}.
Thus the coarse maps are
independent of the choice.
Their group law holds on the
unchanged open set, and then
on the normal chart. This
proves the actual finite action,
including \(g^{d_0}=G^k\)
on a period lift, without
a new linearization cocycle.

More explicitly the new
transverse coordinate is
\(U=p/s^{hm}\). Its finite
map is
\[
             U\longmapsto
       \frac{a_0}{\xi^{hm}}U
                 +\frac{b(s)}{\xi^{hm}s^{hm}} .
 \tag{C.WT.37}\label{cdp:eq:C.WT.37}
\]
On logarithmic fibre coordinates,
the substitution adds \(nh\log s\);
the new off-diagonal coefficient
is \(c(s)s^{hm}\). Both are
holomorphic in the new chart.
Thus subsequent stabilizers
again have the exact affine
form \eqref{cdp:eq:C.WT.28}, with the
remaining displacement in
\eqref{cdp:eq:C.WT.36}. This proves the
induction for the complete
triangular action, not just
its action on homology.

A product ruled cap uses the
actual normalized blowup of
\((p,s^h)\). Its two charts are
\[
                  p=s^hU,\qquad\text{and}\qquad xW=s^h .
 \tag{C.WT.38}\label{cdp:eq:C.WT.38}
\]
The height test has \(\ell>h\);
the central projective normal
action is therefore the identity
for the infinite deck map.
The cap is the specified
product \(F\times\PP^1\), and
its child displacement is
\(\delta(s)/s^h\).
The same finite jet gauge,
now through order \(h-1\),
gives its actual equivariant
blowup; its central affine
map is the recorded finite
projectivized normal map.

Each step subtracts exactly
\(h_em_e\) from the period
order, with \(m_e=1\) for
a product cap. The height
test consequently supplies
the required holomorphic
inequalities at every step
and stops with zero remaining
order only at a nonlinear
terminal leaf. Reverse the
complete group-invariant
pruning record. Equations
\eqref{cdp:eq:C.WT.15}, \eqref{cdp:eq:C.WT.33}--\eqref{cdp:eq:C.WT.38}
restore every original root
component, all its child marks,
and its full torsor and
linear and translation
gluing maps.

\par\medskip\noindent\textbf{Properness of the compressed replacements.}
The affine-base estimate must
retain \(h\). Put
\[
 \rho(p,s)=
       (|p|^2+|s|^{2hm})^{1/(2m)}
       =|x|(|Z|^2+|W|^{2m})^{1/(2m)} .
 \tag{C.WT.39}\label{cdp:eq:C.WT.39}
\]
It satisfies \(\rho\ge|s|^h\).
Since \(\ell\ge hm\), the
translation by \(j\delta(s)\)
distorts this norm by at most
a uniform polynomial in
\(1+|j|\), in both directions.
The pulled line has the same
local fibre \(x^{-n}\) as
before base change. Its norm
distortion is a polynomial
of exponent at most \(n/m\).
Uniform exponential contraction
by \(\gamma(s)^j\) dominates
that polynomial, so a compact
norm annulus meets every orbit
and the action is proper
and free. The finite-group
orbit of a chosen norm gives
the same estimates uniformly.

For comparison with the old
annular chart, the required
normalizing deck exponent is
\(O(1+|\log\rho|)\). Its base
displacement is bounded by
\[
              C|\delta(s)|(1+|\log\rho|).
 \tag{C.WT.40}\label{cdp:eq:C.WT.40}
\]
If \(\rho<1\), use
\(\rho\ge|s|^h\). If \(\rho\ge1\),
an old compact base bound
gives
\(|p_{\rm new}|\le B+
C|\delta(s)|\log(2+|p_{\rm new}|)\),
so \(p_{\rm new}\) is uniformly
bounded. The required new
annular domain is consequently
compact over every old compact
domain, even though the
universal affine base is
not compact.

Finally, on a compact Hopf
fundamental domain, \(x\to0\)
implies \(|s|^h=|xW|=O(|x|)\).
The old base coordinate after
normalization is
\[
          O(|x|^m+|s|^\ell|\log|x||)\longrightarrow0.
 \tag{C.WT.41}\label{cdp:eq:C.WT.41}
\]
Every cluster point is therefore
on the single removed compact
elliptic curve. The closed
overlap graphs and nested
compact annuli used in \ref{cdp:st:C.2.UN.60}
now give the same Hausdorff
proper replacement, retaining
the literal \(A_{h-1}\) neck.
Product caps are proper
normalized blowups. Distinct
central marks have disjoint
finite-group-invariant tubes;
their finite simultaneous
replacements therefore glue.
There are only finitely many
stages.

\par\medskip\noindent\textbf{The smooth weighted quotient and the compact second end.}
The constructed root family
has only its recorded neck
singularities. At a generic
point of a strict component
over \(S_i\), the inertia has
order \(N/m_i\), fixes the
component and acts on its
normal parameter by that
primitive character. Its
quotient is smooth and the
full fibre multiplicity is
exactly \(m_i\).

At an original double curve
of multiplicities \(a=gA,b=gB\),
the stabilizer of a point of
one lifted curve has order
\(hAB=N/g\). Its subgroup
\(\mu_h\) fixes \(u,v\) and
acts primitively on \(s\).
Quotienting \(uv=s^h\) first
by it gives the smooth
\((u,v)\)-chart. The residual
cyclic group of order \(AB\)
has the two independent
characters of orders \(B,A\).
Since \(\gcd(A,B)=1\), its
invariant ring is exactly
\[
 \bigl(\C\{z,u,v,s\}/(uv-s^h)\bigr)^{\,I}
       =\C\{z,u^B,v^A\},\qquad
                  t=(u^B)^a(v^A)^b .
 \tag{C.WT.42}\label{cdp:eq:C.WT.42}
\]
Here \(I\) is that inertia
group. The curve coordinate
\(z\) is invariant, because
the actual curve map is
\'{e}tale. Finite character
averaging gives coordinates
of the form \(uU,vV,z+\cdots\)
with their original characters.
Their product is \(s^h\)
times an invariant unit;
dividing one coordinate by
that unit keeps both the
characters and the literal
equation \(uv=s^h\).
This proves \eqref{cdp:eq:C.WT.42} for the
actual finite action, not
only its tangent representation.

There is no edge inversion
to add an unrecorded quotient
of this chart. At interior
points the strict component
maps are \'{e}tale after the
generic normal inertia, so
there are no other quotient
singularities. The resulting
total space near zero is
smooth and its weighted
central support is the
original one.

For completeness the weighted
normal line is retained as
well. On the root complement
the original normal coordinate
is \(v_i=s^{N/m_i}/\xi\).
The character eigensheaves
of the finite component map
extend by \eqref{cdp:eq:C.WT.7}; explicitly
their \(k\)-th pieces are
\[
 \mathcal L_i^{\,k}\otimes
  \OO_{S_i}\!\left(
       \sum_j\left\lfloor\frac{k m_j}{m_i}\right\rfloor J_{ij}
             \right),\qquad 0\le k<m_i .
 \tag{C.WT.43}\label{cdp:eq:C.WT.43}
\]
This follows by dividing
\(\xi^k\) by its integral
normal zero order in that
chart. The retained full
section and its orders
recover \(\mathcal L_i\)
and the specified edge line
maps; no reduction of an
integer modulo \(m_i\)
has occurred. The chosen
normal-product constants
are the only freedom stated
in the theorem.

Carry out the conjugate
pruning record at infinity
as well. Formula \eqref{cdp:eq:C.WT.25}
has base character
\(\zeta_N^{-1}\) and canonical
character \((\det D)^{-1}\);
its primitive elliptic period
again has exact order \(R\).
Thus the conjugate height
and finite-jet comparisons
hold with precisely the
conjugate characters.
All charts, including the
compressed necks and
\eqref{cdp:eq:C.WT.42}, are the actual
conjugates of those at zero.
The quotient is therefore
smooth at that end too and
has the conjugate prescribed
weighted support. The two
sets of replacements have
disjoint parameter
neighbourhoods.

The compact core and the
proper replacements give
a compact normal root
threefold. Its finite
quotient is smooth by the
charts just proved. Away
from zero and infinity
the base action is free;
one fibre over each orbit
maps biholomorphically to
the quotient fibre. They
are consequently connected
complex two-tori. Each
original component and
each double elliptic curve
has Euler number zero,
so both reduced central
supports have Euler zero.
This completes necessity,
sufficiency and the compact
comparison, including all
finite triangular actions.
\end{proof}


